%------------------------------------------------------------------------------%
\begin{question}
Determine a forma escalonada reduzida da matriz
\[
  A=
  \begin{bmatrix}
    1  & 2 & -1 \\
    2  & 1 & 1  \\
    -1 & 3 & -2
  \end{bmatrix}
\]
\end{question}

%------------------------------------------------------------------------------%
\begin{resolution}{Escalonamento}
\begin{columns}
\column{0.3\linewidth}
  \[\;\;\;
  \begin{bmatrix}
    1  & 2 & -1 \\
    2  & 1 & 1  \\
    -1 & 3 & -2
  \end{bmatrix}
  \]
  \pause
\column{0.7\linewidth}
  Eliminamos os elementos abaixo do primeiro pivô
  \\ \pause
  \(
    L_2 \leftarrow L_2 -2 L_1
  \)
  \\ \pause
  \(
    L_3 \leftarrow L_3 + L_1
  \)
\end{columns}
\vspace{-\baselineskip}
\[
  \pause
  L'_2
  \;=\;
  L_2 -2 L_1
  \pause
  \;=\;
  \begin{bmatrix}
    2  & 1 & 1
  \end{bmatrix}
  -2
  \begin{bmatrix}
    1  & 2 & -1
  \end{bmatrix}
  \pause
  \;=\;
  \begin{bmatrix}
    0 & -3 & 3
  \end{bmatrix}
\]
\[
  \pause
  L'_3
  \;=\;
  L_3 + L_1
  \pause
  \;=\;
  \begin{bmatrix}
    -1 & 3 & -2
  \end{bmatrix}
  +
  \begin{bmatrix}
    1  & 2 & -1
  \end{bmatrix}
  \pause
  \;=\;
  \begin{bmatrix}
    0 & 5  & -3
  \end{bmatrix}
\]
\pause
\[
  \begin{bmatrix}
    1 & 2  & -1 \\
    0 & -3 & 3  \\
    0 & 5  & -3
  \end{bmatrix}
\]
\end{resolution}

%------------------------------------------------------------------------------%
\begin{resolution}{Escalonamento}
\begin{columns}
\column{0.3\linewidth}
  \[\;\;\;
  \begin{bmatrix}
    1 & 2  & -1 \\
    0 & -3 & 3  \\
    0 & 5  & -3
  \end{bmatrix}
  \]
  \pause
\column{0.7\linewidth}
  Fazendo com que o segundo pivô seja 1
  \\ \pause
  \(
  L_2 \leftarrow -\frac{1}{3}L_2
  \)
\end{columns}
\vspace{-\baselineskip}
\[
  \pause
  L'_2
  \;=\;
  \pause
  \;=\;
  -\frac{1}{3}L_2
  \begin{bmatrix}
    0 & -3 & 3
  \end{bmatrix}
  \pause
  \;=\;
  \begin{bmatrix}
    0 & 1 & -1
  \end{bmatrix}
\]
\pause
\[
  \begin{bmatrix}
    1 & 2 & -1 \\
    0 & 1 & -1 \\
    0 & 5 & -3
  \end{bmatrix}
\]
\end{resolution}

%------------------------------------------------------------------------------%
\begin{resolution}{Escalonamento}
\begin{columns}
\column{0.3\linewidth}
  \[\;\;\;
  \begin{bmatrix}
    1 & 2 & -1 \\
    0 & 1 & -1 \\
    0 & 5 & -3
  \end{bmatrix}
  \]
  \pause
\column{0.7\linewidth}
  Eliminamos os elementos abaixo do segundo pivô
  \\ \pause
  \(
    L_3 \leftarrow L_3 - 5 L_2
  \)
\end{columns}
\vspace{-\baselineskip}
\[
  \pause
  L'_3
  \;=\;
  L_3 - 5 L_2
  \pause
  \;=\;
  \begin{bmatrix}
    0 & 5 & -3
  \end{bmatrix}
  - 5
  \begin{bmatrix}
    0 & 1 & -1
  \end{bmatrix}
  \pause
  \;=\;
  \begin{bmatrix}
    0 & 0 & 2
  \end{bmatrix}
\]
\pause
\[
  \begin{bmatrix}
    1 & 2 & -1 \\
    0 & 1 & -1 \\
    0 & 0 & 2
  \end{bmatrix}
\]
\end{resolution}

%------------------------------------------------------------------------------%
\begin{resolution}{Escalonamento}
\begin{columns}
\column{0.3\linewidth}
  \[\;\;\;
  \begin{bmatrix}
    1 & 2 & -1 \\
    0 & 1 & -1 \\
    0 & 0 & 2
  \end{bmatrix}
  \]
  \pause
\column{0.7\linewidth}
  Fazendo com que o terceiro pivô seja 1
  \\ \pause
  \(
    L_3 \leftarrow \frac{1}{2} L_3
  \)
\end{columns}
\vspace{-\baselineskip}
\[
  \pause
  L'_3
  \;=\;
  \frac{1}{2} L_3
  \pause
  \;=\;
  \frac{1}{2}
  \begin{bmatrix}
    0 & 0 & 2
  \end{bmatrix}
  \pause
  \;=\;
  \begin{bmatrix}
    0 & 0 & 1
  \end{bmatrix}
\]
\pause
\[
  \begin{bmatrix}
    1 & 2 & -1 \\
    0 & 1 & -1 \\
    0 & 0 & 1
  \end{bmatrix}
\]
\end{resolution}

%------------------------------------------------------------------------------%
\begin{resolution}{Escalonamento}
\begin{columns}
\column{0.3\linewidth}
  \[\;\;\;
  \begin{bmatrix}
    1 & 2 & -1 \\
    0 & 1 & -1 \\
    0 & 0 & 1
  \end{bmatrix}
  \]
  \pause
\column{0.7\linewidth}
  Eliminamos os elementos acima do terceiro pivô
  \\ \pause
  \(
    L_2 \leftarrow L_2 + L_3
  \)
  \\ \pause
  \(
    L_1 \leftarrow L_1 + L_3
  \)
\end{columns}
\vspace{-\baselineskip}
\[
  \pause
  L'_2
  \;=\;
  L_2 + L_3
  \pause
  \;=\;
  \begin{bmatrix}
    0 & 1 & -1
  \end{bmatrix}
  +
  \begin{bmatrix}
    0 & 0 & 1
  \end{bmatrix}
  \pause
  \;=\;
  \begin{bmatrix}
    0 & 1 & 0
  \end{bmatrix}
\]
\[
  \pause
  L'_1
  \;=\;
  L_1 + L_3
  \pause
  \;=\;
  \begin{bmatrix}
    1 & 2 & -1
  \end{bmatrix}
  +
  \begin{bmatrix}
    0 & 0 & 1
  \end{bmatrix}
  \pause
  \;=\;
  \begin{bmatrix}
    1 & 2 & 0
  \end{bmatrix}
\]
\pause
\[
  \begin{bmatrix}
    1 & 2 & 0 \\
    0 & 1 & 0 \\
    0 & 0 & 1
  \end{bmatrix}
\]
\end{resolution}

%------------------------------------------------------------------------------%
\begin{resolution}{Escalonamento}
\begin{columns}
\column{0.3\linewidth}
  \[\;\;\;
  \begin{bmatrix}
    1 & 2 & 0 \\
    0 & 1 & 0 \\
    0 & 0 & 1
  \end{bmatrix}
  \]
  \pause
\column{0.7\linewidth}
  Eliminamos os elementos acima do segundo pivô
  \\ \pause
  \(
    L_1 \leftarrow L_1 - 2 L_2
  \)
\end{columns}
\vspace{-\baselineskip}
\[
  \pause
  L'_1
  \;=\;
  L_1 - 2 L_2
  \pause
  \;=\;
  \begin{bmatrix}
    1 & 2 & 0
  \end{bmatrix}
  - 2
  \begin{bmatrix}
    0 & 1 & 0
  \end{bmatrix}
  \pause
  \;=\;
  \begin{bmatrix}
    1 & 0 & 0
  \end{bmatrix}
\]
\pause
\[
  \begin{bmatrix}
    1 & 0 & 0 \\
    0 & 1 & 0 \\
    0 & 0 & 1
  \end{bmatrix}
\]

\pause
Forma escalonada reduzida da matriz $A$
\end{resolution}

%------------------------------------------------------------------------------%
\endinput
